\section{Simultaneous normalization of all rational Kummer coordinates}
\label{sec:rational-radical-normalization}

The arithmetic part of every rational Kummer coordinate can be removed
at once. This replaces the normalized subgroup by a proper closed
subgroup without changing the relative abelian obstruction.
The Galois calculation used here is the classical description of radical
extensions; see \cite[\S2, especially Theorem~2.5, author pp.~4--6]{MoreeStevenhagen}.
We give the finite-coordinate argument to make its simultaneous use and
the resulting GT quotient explicit.

\begin{lemma}\label{lem:simultaneous-kummer-image}
For distinct rational primes $q_1,\ldots,q_m$, the joint Kummer map
\[
 G_{\mathbb Q^{\mathrm{ab}}}\longrightarrow\widehat{\mathbb Z}^{\,m},
 \qquad \sigma\longmapsto
       (\kappa_{q_1}(\sigma),\ldots,\kappa_{q_m}(\sigma))
\]
has image exactly $2\widehat{\mathbb Z}^{\,m}$.
Consequently, for distinct odd primes $q_1,\ldots,q_m$, its restriction
to $K=G_{K_2}$ also has image $2\widehat{\mathbb Z}^{\,m}$.
\end{lemma}

\begin{proof}
Denote the first image by $J$. It is a compact, hence closed, subgroup
of $\widehat{\mathbb Z}^{\,m}$. Every continuous character of that
product with values in $\mathbb Q/\mathbb Z$ has the form
\[
 (t_1,\ldots,t_m)\longmapsto\frac1n\sum_{i=1}^m a_it_i\pmod{\mathbb Z}
\]
for integers $n>0$ and $a_i$. Such a character vanishes on $J$ exactly
when the positive real radical
\[
 \alpha=\prod_{i=1}^m q_i^{a_i/n}
\]
belongs to $\mathbb Q^{\mathrm{ab}}$. Indeed the Kummer action on
$\alpha$ multiplies it by the corresponding $n$th root of unity.

If $\alpha\in\mathbb Q^{\mathrm{ab}}$, its field over $\mathbb Q$
is normal and is contained in $\mathbb R$. Every conjugate of
$\alpha$ is a real root of $X^n-\alpha^n$, and is therefore one
of $\alpha,-\alpha$. It follows that $\alpha^2\in\mathbb Q$.
Writing this last rational number in its unique prime factorization
gives $2a_i/n\in\mathbb Z$ for each $i$. Conversely these
integrality conditions make $\alpha$ a square root of a positive
rational number, and hence an element of $\mathbb Q^{\mathrm{ab}}$.
Thus the annihilator of $J$ is precisely the group of characters
of exponent at most two. It is also the annihilator of
$2\widehat{\mathbb Z}^{\,m}$, so the two closed subgroups agree.
For this last implication, a proper closed subgroup of a profinite
abelian group is separated from any point outside it by a character
of a finite quotient.

For the second assertion, apply the first to the list
$2,q_1,\ldots,q_m$. Its image is the full product
$2\widehat{\mathbb Z}^{\,m+1}$. Every point of that product whose
first coordinate is exactly zero has a preimage fixing all roots
of two as well as all roots of unity. Such a preimage belongs to
$K$, proving the assertion with no compatibility condition left over.
\end{proof}

Let $\mathcal P$ be the set of rational primes and put
\[
 L_{\mathrm{rad}}=
 \mathbb Q\bigl(\mu_\infty,q^{1/n}:q\in\mathcal P,
                                      \ n\geq1\bigr).
\]
The roots in this expression are positive real roots; adjoining
$\mu_\infty$ also adjoins all their conjugates.
Define
\[
 N_{\mathrm{rad}}=N\cap
       \bigcap_{\substack{q\in\mathcal P\\q\ne2}}\ker\rho_q^I,
 \qquad K_{\mathrm{rad}}=G_{L_{\mathrm{rad}}}.
\]
The frame cocycle $\rho_2$ and $\rho_2^I$ agree on $G$;
the normalization is recorded in \cite[\S2.2, p.~100, footnote]{Furusho2003}
and \cite[Proposition~11]{H17}. Thus $N_{\mathrm{rad}}$ is the
simultaneous kernel of the multiplier and all rational-prime Kummer
cocycles.

\begin{theorem}\label{thm:rational-radical-normalization}
The joint map
\[
 \mathcal R:N\longrightarrow
       \prod_{\substack{q\in\mathcal P\\q\ne2}}2\widehat{\mathbb Z},
 \qquad n\longmapsto(\rho_q^I(n))_q
\]
is surjective already on $K$. Its kernel is $N_{\mathrm{rad}}$,
and
\[
 N=N_{\mathrm{rad}}K,\qquad
 G=N_{\mathrm{rad}}H,\qquad
 N_{\mathrm{rad}}\cap H=K_{\mathrm{rad}}.
\]
In particular $N_{\mathrm{rad}}$ is a proper closed normal subgroup
of $N$, and
\[
 N/N_{\mathrm{rad}}\simeq
       \prod_{\substack{q\in\mathcal P\\q\ne2}}2\widehat{\mathbb Z}.
\]
There are canonical topological isomorphisms
\begin{align*}
 G/N_{\mathrm{rad}}&\simeq
            \operatorname{Gal}(L_{\mathrm{rad}}/\mathbb Q),\\
 B&\simeq N_{\mathrm{rad}}/
       \overline{[G,N_{\mathrm{rad}}]K_{\mathrm{rad}}}.
\end{align*}
The first quotient is explicitly
\[
 A_{\mathrm{rad}}=
 \left\{(u,(t_q)_{q\in\mathcal P}):
 \begin{array}{l}
 u\in\widehat{\mathbb Z}^{\times},\quad
 t_q\in\widehat{\mathbb Z},\\
 t_q\equiv\epsilon_q(u)\pmod2\quad(q\in\mathcal P)
 \end{array}\right\},
\]
with multiplication $(u,t)(v,s)=(uv,t+us)$, where
$\epsilon_q$ is the quadratic character of $\mathbb Q(\sqrt q)$.
\end{theorem}

\begin{proof}
Each coordinate of $\mathcal R$ is a homomorphism on $N$ and is
even by Theorem~\ref{thm:all-kummer-parity}. The joint map is
continuous for the product topology. By
Lemma~\ref{lem:simultaneous-kummer-image}, every finite-coordinate
projection of $\mathcal R(K)$ is the whole corresponding product.
Hence $\mathcal R(K)$ is dense in the full product. It is compact
and therefore closed, proving
\[
 \mathcal R(K)=\mathcal R(N)=\prod_{q\ne2}2\widehat{\mathbb Z}.
\]
Equivalently, one can prescribe successively larger finite lists
of coordinates and use the finite intersection property in $K$.
This compactness step prescribes the whole profinite value of every
coordinate, not only its residue at one finite level.

Given $n\in N$, choose $k\in K$ with $\mathcal R(k)=\mathcal R(n)$.
Then $nk^{-1}\in N_{\mathrm{rad}}$, proving
$N=N_{\mathrm{rad}}K$ and then $G=N_{\mathrm{rad}}H$.
The conjugation identity for the Kummer cocycles shows that
$N_{\mathrm{rad}}$ is normal in $G$; it is closed by definition.
Its Galois intersection fixes precisely the roots of unity and
the roots of all rational primes, so it is $K_{\mathrm{rad}}$.
The quotient isomorphism $N/N_{\mathrm{rad}}\simeq\prod_{q\ne2}
2\widehat{\mathbb Z}$ follows from compactness. This nontrivial
quotient also proves that $N_{\mathrm{rad}}$ is proper.

The factorization $G=N_{\mathrm{rad}}H$ identifies
$G/N_{\mathrm{rad}}$ with $H/K_{\mathrm{rad}}$.
On the Galois group, each translation coordinate has parity
$\epsilon_q(u)$ because it records the action on $\sqrt q$.
For a fixed multiplier choose one Galois lift. The first assertion
of Lemma~\ref{lem:simultaneous-kummer-image}, followed by compactness
over all finite prime lists, allows an arbitrary even translation
vector to be added to this lift. Thus the image is exactly
$A_{\mathrm{rad}}$, with the displayed cocycle multiplication.

Finally apply the proof of Theorem~\ref{thm:abelian-splitting}
to $G=N_{\mathrm{rad}}H$. The Galois intersection
$K_{\mathrm{rad}}$ lies in $\overline{[H,H]}$, since it fixes
$\mathbb Q^{\mathrm{ab}}$. That proof identifies the new relative
abelian quotient with
$(G/\overline{\langle\!\langle H\rangle\!\rangle})^{\mathrm{ab}}$,
which is also $B$. This gives the asserted canonical isomorphism.
\end{proof}

\begin{corollary}\label{cor:rational-kummer-characters}
The invariant quadratic characters of $N$ that factor through all
rational Kummer coordinates form exactly
\[
 X_2^{\mathrm{Kum}}=
 \bigoplus_{\substack{q\in\mathcal P\\q\ne2}}
       \mathbb F_2\,\psi_q^N.
\]
Restriction to $K$ is injective on this subspace. For an odd prime
$p$, no nonzero invariant $\mathbb F_p$-character of $N$ factors
through the rational Kummer coordinates.
For every prime $p$, the unresolved relative kernel is equivalently
\[
 \operatorname{Hom}(B,\mathbb F_p)=
 \left\{\chi\in H^1(N_{\mathrm{rad}},\mathbb F_p)^G:
                       \chi|_{K_{\mathrm{rad}}}=0\right\}.
\]
\end{corollary}

\begin{proof}
Enriquez's distribution theorem gives
$\rho_{ab}^I=\rho_a^I+\rho_b^I$ on $G$ for positive integers
$a,b$; see \cite[Theorem~8.1 and Remark~8.6, author pp.~30--32]{Enriquez2007}.
Since $\rho_2^I$ is zero on $N$, all rational Kummer coordinates
on $N$ are generated by the odd-prime coordinates.
A continuous character of the product in
Theorem~\ref{thm:rational-radical-normalization} to a finite group
depends on finitely many coordinates: its open kernel contains
a basic product neighborhood of zero. Each factor
$2\widehat{\mathbb Z}$ has a one-dimensional quadratic character
space, generated by division by two followed by reduction modulo
two. This proves the direct-sum formula. Every multiplier is odd,
so every one of these characters is $G$-invariant. Surjectivity
already on $K$ proves injectivity of restriction.

Complex conjugation acts on the product by multiplication by $-1$.
If $p$ is odd and a character $\chi$ is invariant, then
$\chi(t)=\chi(-t)=-\chi(t)$, so it is zero. The final identification
is the character-gluing lemma applied to
$G=N_{\mathrm{rad}}H$.
\end{proof}

\begin{proposition}\label{prop:radical-character-pullback}
For every prime $p$, put
\[
 X_p^{\mathrm{rad}}=H^1(N_{\mathrm{rad}},\mathbb F_p)^G,
 \qquad
 Y_p^{\mathrm{rad}}=H^1(K_{\mathrm{rad}},\mathbb F_p)^H.
\]
The restriction maps identify $X_p$ with the fiber product
\[
 X_p\simeq X_p^{\mathrm{rad}}
                  \mathop{\times}_{Y_p^{\mathrm{rad}}}Y_p.
\]
The kernels of $X_p\to X_p^{\mathrm{rad}}$ and
$Y_p\to Y_p^{\mathrm{rad}}$ are identified by restriction.
They are zero for odd $p$, and for $p=2$ they are respectively
$X_2^{\mathrm{Kum}}$ and its restriction to $K$.
\end{proposition}
\begin{proof}
A character of $N$ restricts to the indicated compatible pair,
and the pair determines it because $N=N_{\mathrm{rad}}K$.
Conversely, let $\alpha\in X_p^{\mathrm{rad}}$ and
$\beta\in Y_p$ have the same restriction to $K_{\mathrm{rad}}$.
Define
\[
 \chi(nk)=\alpha(n)+\beta(k)
       \qquad(n\in N_{\mathrm{rad}},\ k\in K).
\]
If $nk=n'k'$, then $n'^{-1}n=k'k^{-1}$ lies in
$N_{\mathrm{rad}}\cap K=K_{\mathrm{rad}}$; the equality of
the two restrictions proves well-definedness. The invariance of
$\alpha$ under $K$ makes this formula a homomorphism.
It is continuous: the multiplication map from the compact space
$N_{\mathrm{rad}}\times K$ onto the Hausdorff group $N$ is
a quotient map, and its pullback is the continuous sum
$\alpha+\beta$. Conjugation by $H$ preserves the formula,
since $\alpha$ is $G$-invariant and $\beta$ is $H$-invariant.
Conjugation by $N$ preserves any abelian character of $N$.
Since $G=NH$, the resulting character belongs to $X_p$.

The kernel of restriction to $N_{\mathrm{rad}}$ is the invariant
character space of $N/N_{\mathrm{rad}}$; the corresponding kernel
on the Galois side is that of $K/K_{\mathrm{rad}}$.
The joint Kummer map identifies both quotients with
$\prod_{q\ne2}2\widehat{\mathbb Z}$, with the same scalar action
of the multiplier. The preceding corollary computes their
invariant character spaces, proving the final assertions.
\end{proof}

The subgroup $N_{\mathrm{rad}}$ gives a strictly smaller location
for the open obstruction. A character of $G$ which vanishes on
$H$ is zero if and only if its restriction to this subgroup is zero.
The theorem does not determine its relative abelianization.
